% =====================================================================
%  CONTENUTO DELLA VERIFICA — qui si scrivono dati ed esercizi.
%  Per ottenere i PDF imposta come Main document uno dei file numerati
%  di questa cartella (1_Testo_FilaA, 4_Soluzioni_FilaA, 7_Correzione, ...).
%  Nuova verifica: copia la cartella come
%     Compiti/Verifica_2026_10_27_NomeMateria_3A/
%  e aggiorna il percorso nell'ultima riga dei file numerati.
% =====================================================================
% valori predefiniti se compili direttamente questo file
\providecommand{\fila}{A}          % A | B | C
\providecommand{\versione}{testo}  % testo | soluzioni | correzione | completa

\documentclass[a4paper,11pt,addpoints]{exam}
\usepackage{stili/verifica}

% ---------- Dati della prova (lascia vuoto ciò che non serve) ----------
\scuola{}               % es. Istituto ..., Città
\annoscolastico{}       % es. 2026/27  (vuoto -> riga da compilare)
\tipoprova{Verifica}    % Verifica, Compito in classe, Prova di recupero...
\materia{NomeMateria}   % titolo: "Verifica di NomeMateria" (vuoto -> solo "Verifica")
\argomento{}            % sottotitolo
\classe{}               % es. 3ª  (vuoto -> riga da compilare)
\dataprova{}            % vuoto -> riga da compilare
\durata{}               % es. 60 minuti
\votominimo{2}          % voto con 0 punti
\sogliasufficienza{60}  % % del punteggio per il 6
\criterigenerali{Risultato senza unità di misura o con unità errata: $-0{,}5$ pt.
  Errore di calcolo con procedimento corretto: metà dei punti della richiesta.
  Risultati accettati con 2--3 cifre significative.}   % nella griglia di correzione

\begin{document}

\intestazione

\begin{questions}

% Esercizi presi dalla banca (un file = un esercizio, percorso dalla radice)
% ---------------------------------------------------------------------
% Argomento: resistenze serie/parallelo, legge di Ohm, potenza
% Livello: base-intermedio   |   Tempo stimato: 12 min
% ---------------------------------------------------------------------
\begin{esercizio}
\dato{V}{12}{15}{18}
\dato{R1}{150}{180}{220}
\dato{R2}{220}{270}{330}
\dato{R3}{330}{560}{680}
\calcola[R_{23}]{R23}{\x{R2}*\x{R3}/(\x{R2}+\x{R3})}
\calcola[R_{eq}]{Req}{\x{R1}+\x{R23}}
\calcola[I]{I}{\x{V}/\x{Req}}
\calcola[P]{P}{\x{V}*\x{I}}

Dato il circuito in figura:

\begin{center}
\begin{minipage}{0.48\linewidth}\centering
  \begin{circuitikz}[scale=0.85, transform shape]
    \draw (0,0) to[vsource, l=$V$, n=G] (0,3)
          to[R, l=$R_1$, -*] (3,3) -- (5,3)
          (3,3) to[R, l_=$R_2$, -*] (3,0)
          (5,3) to[R, l=$R_3$] (5,0) -- (0,0);
    \morsetti{G}
  \end{circuitikz}
\end{minipage}\hfill
\begin{minipage}{0.4\linewidth}
  \renewcommand{\arraystretch}{1.3}
  \begin{tabular}{r@{${}={}$}l}
    $V$   & \mis{V}{\volt}\\
    $R_1$ & \mis{R1}{\ohm}\\
    $R_2$ & \mis{R2}{\ohm}\\
    $R_3$ & \mis{R3}{\ohm}
  \end{tabular}
\end{minipage}
\end{center}

\begin{parts}
  \part[2] Calcola la resistenza equivalente $R_{eq}$ vista dal generatore.
  \part[2] Calcola la corrente $I$ erogata dal generatore.
  \part[2] Calcola la potenza $P$ erogata dal generatore.
\end{parts}

\begin{solution}
  $R_{23} = \dfrac{R_2 R_3}{R_2+R_3} = \dfrac{\val{R2}\cdot\val{R3}}{\val{R2}+\val{R3}}
    \approx \mis{R23}{\ohm}$,\qquad
  $R_{eq} = R_1 + R_{23} \approx \mis{Req}{\ohm}$\\[4pt]
  $I = V/R_{eq} \approx \mis{I}{\ampere}$\qquad
  $P = V\cdot I \approx \mis{P}{\watt}$
  \criterio{a) 1 pt riconoscere il parallelo $R_2\parallel R_3$, 1 pt calcolo;
    b) e c) 1 pt formula, 1 pt risultato con unità.}
\end{solution}
\end{esercizio}

% ---------------------------------------------------------------------
% Argomento: legge di Ohm, potenza dissipata
% Livello: base   |   Tempo stimato: 6 min
% ---------------------------------------------------------------------
\begin{esercizio}
\dato{R}{470}{680}{560}
\dato{V}{9}{12}{15}
\calcola[I]{I}{\x{V}/\x{R}}
\calcola[P]{P}{\x{V}^2/\x{R}}

Un resistore da \mis{R}{\ohm} è collegato a una tensione di \mis{V}{\volt}.
\begin{parts}
  \part[2] Calcola la corrente che lo attraversa.
  \part[2] Calcola la potenza dissipata.
\end{parts}

\begin{solution}
  $I = V/R \approx \mis{I}{\ampere}$;\qquad
  $P = V^2/R \approx \mis{P}{\watt}$
  \criterio{1 pt formula, 1 pt risultato; $-0{,}5$ se manca o è sbagliata l'unità.}
\end{solution}
\end{esercizio}


% ---------------------------------------------------------------------
% Argomento: resistenze in parallelo
% Livello: intermedio   |   Tempo stimato: 6 min
% ---------------------------------------------------------------------
\begin{esercizio}[4]
\dato{Ra}{1e3}{1.5e3}{1.2e3}
\dato{Rb}{2.2e3}{3.3e3}{1.8e3}
\dato{Rc}{3.3e3}{4.7e3}{2.7e3}
\calcola[R_{eq}]{Req}{1/(1/\x{Ra}+1/\x{Rb}+1/\x{Rc})}

Tre resistori da \mis{Ra}{\ohm}, \mis{Rb}{\ohm} e \mis{Rc}{\ohm}
sono collegati in parallelo. Calcola la resistenza equivalente.

\begin{solution}
  $\dfrac{1}{R_{eq}} = \dfrac{1}{R_a}+\dfrac{1}{R_b}+\dfrac{1}{R_c}
   \;\Rightarrow\; R_{eq} \approx \mis{Req}{\ohm}$
  \criterio{2 pt formula del parallelo, 2 pt calcolo; accettare
    arrotondamenti a 2--3 cifre.}
\end{solution}
\end{esercizio}


% Oppure un esercizio scritto direttamente qui (vedi Banca_Esercizi/_MODELLO_esercizio.tex)

\end{questions}

% Tabella punti + griglia voti in fondo alla 2ª pagina
\riepilogo

\end{document}
